\section{Discussion}

The evaluation shows why image-manipulation detector performance should be read jointly across datasets, thresholds and transformations. Both frozen models achieved their highest AUC on Columbia Uncompressed, while their AUCs were lower on corrected CASIA v1 and COVERAGE. Under the fixed 0.5 threshold, ranking quality and binary operating performance were not interchangeable: MVSS-Net's Columbia AUC was high, but its false-positive count on authentic images was also high. A single AUC therefore does not describe the decision behavior relevant to a specified operating threshold.

Localization results were also conditional on both manipulation type and reference-mask semantics. On CASIA v1, both models produced higher macro F1 for splicing than for copy-move manipulation, with a larger difference for MVSS-Net. On Columbia, the provided blue-region reference denotes a dataset-specific boundary/source-region convention rather than the complete edited object. Comparing an anomaly map to that reference measures spatial agreement with the provided annotation, not necessarily recovery of the full manipulated region. The dimension mismatch in nine COVERAGE masks similarly required those cases to be excluded from localization while retaining them for image-level detection. These distinctions should accompany any interpretation of the aggregate localization scores.

The controlled stress tests demonstrate sensitivity to post-processing in the tested samples. JPEG recompression and resize--restore affected the two datasets and models differently; strong resizing caused particularly large degradation for TruFor on CASIA v1. Since thresholds were fixed, these results expose changes in both score ordering and operating-point behavior. They do not establish the behavior of all detectors, social-media platforms, screen recapture, unseen cameras or future image-processing pipelines. The transformations are deterministic digital stress conditions, not a complete model of deployment distribution shift.

The evidence audit contributes a separate computational finding. Complete evidence records and successful independent-process replay show that the sampled predictions can be regenerated from hash-matched inputs and frozen artifacts under the recorded environment. Replay agreement does not imply prediction correctness, causal attribution, or diagnostic uniqueness. Nor does an independent process on the same host demonstrate cross-platform, cross-machine or third-party-laboratory reproducibility. No practitioner study was conducted, so no claim is made about improved human judgment, review time or report quality.

Several limits define the next validation steps. Only two pretrained models and three existing datasets were evaluated, and the potential overlap between the CASIAv2 training source and the corrected CASIA v1 bundle qualifies that comparison. The copy-move external evaluation in COVERAGE contains only 200 image-level cases, with 91 usable manipulation masks. Columbia's established status also means its high score should be interpreted as a known-benchmark replication rather than a wholly unseen-domain result. Future work should include additional independently sourced datasets, model versions, cross-machine replay, stronger split audits, and a preregistered examiner study before making broader operational claims. A small dual-branch training pilot is reported as supplementary negative evidence; it did not demonstrate a reliable external-domain gain and is not used to claim algorithmic novelty.
